% Agent 01: Pages 345--348 (Title, Table of Contents, Section 1, Section 2.1 start)
% =============================================================================

% --- Title Page (p.~345) ---

\title{Astrophysics of Black Holes}

\author{Igor~D.~Novikov}
\thanks{Supported in part by the Academy of Sciences of the USSR}
\affiliation{Institute of Applied Mathematics, Academy of Sciences, Moscow}

\author{Kip~S.~Thorne}
\thanks{Supported in part by the U.S.\ National Science Foundation [GP-27304, GP-28027]}
\affiliation{California Institute of Technology, Pasadena}

\maketitle

% --- Table of Contents (pp.~345--346) ---

\tableofcontents

% NOTE: The original Table of Contents lists the following sections and page numbers:
%
% 1   Introductory Remarks ...................................... 347
% 2.1 Thermal Bremsstrahlung (``Free-Free Radiation'' From a Plasma) 347
% 2.2 Free-Bound Radiation ...................................... 355
% 2.3 Thermal Cyclotron and Synchrotron Radiation ............... 357
% 2.4 Electron Scattering of Radiation .......................... 360
% 2.5 Hydrodynamics and Thermodynamics .......................... 362
% 2.6 Radiative Transfer ........................................ 370
% 2.7 Shock Waves ............................................... 379
% 2.8 Turbulence ................................................ 383
% 2.9 Reconnection of Magnetic Field Lines ...................... 384
% 3   The Origin of Stellar Black Holes ......................... 386
% 4   Black Holes in The Interstellar Medium ..................... 389
% 4.1 Accretion of Noninteracting Particles onto a Nonmoving
%     Black Hole ................................................. 389
% 4.2 Adiabatic, Hydrodynamic Accretion onto a Nonmoving
%     Black Hole ................................................. 390
% 4.3 Thermal Bremsstrahlung from the Accreting Gas ............. 396
% 4.4 Influence of Magnetic Fields and Synchrotron Radiation ..... 397
% 4.5 Interaction of Outflowing Radiation with the Gas ........... 399
% 4.6 Validity of the Hydrodynamical Approximation ............... 401
% 4.7 Gas Flow Near the Horizon .................................. 402
% 4.8 Accretion onto a Moving Hole ............................... 404
% 4.9 Optical Appearance of Hole: Summary ........................ 407
% 5   Black Holes in Binary Star Systems and in the Nuclei
%     of Galaxies ................................................ 408
% 5.1 Introduction ............................................... 408
% 5.2 Accretion in Binary Systems: The General Picture ........... 409
% 5.3 Accretion in Galactic Nuclei: The General Picture .......... 420
% 5.4 Properties of the Kerr Metric Relevant to Accreting Disks .. 422
% 5.5 Relativistic Model for Disk: Underlying Assumptions ........ 424
% 5.6 Equations of Radial Structure .............................. 427
% 5.7 Equations of Vertical Structure ............................ 431
% 5.8 Approximate Version of Vertical Structure .................. 434
% 5.9 Explicit Models for Disk .................................. 435
% 5.10 Spectrum of Radiation from Disk ........................... 438
% 5.11 Heating of the Outer Region by X-Rays from the
%      Inner Region .............................................. 443
% 5.12 Fluctuations on the Steady State Model .................... 443
% 5.13 Supercritical Accretion .................................. 444
% 5.14 Comparison with Observations ............................. 445
% 6   White Holes and Black Holes of Cosmological Origin ......... 445
% 6.1 White Holes, Grey Holes, and Black Holes ................... 445
% 6.2 The Growth of Cosmological Holes by Accretion .............. 446
% 6.3 Limit on the Number of Baryons in Cosmological Holes ....... 447
% 6.4 Caution .................................................... 448
% References ..................................................... 448


% =============================================================================
% SECTION 1: Introductory Remarks (p.~347)
% =============================================================================

\section{Introductory Remarks}
\label{sec:1}

To analyze astrophysical aspects of black holes one must bring input from two
directions: from the elegant realm of general relativity (lectures of Hawking,
Carter, Bardeen, and Ruffini), and from the more mundane realm of astrophysics
and plasma physics (\S\ref{sec:2} of these lectures).  We shall here combine these
inputs to discuss the origin of stellar black holes (\S\ref{sec:3}); and to analyze observable
aspects of black holes in the interstellar medium (\S\ref{sec:4}), in binary star systems and
the nuclei of galaxies (\S\ref{sec:5}), and in cosmological contexts (\S\ref{sec:6}).

Readers with strong astrophysical backgrounds will wish to skip over our
brief treatment of the fundamentals of astrophysics and plasma physics (\S\ref{sec:2}),
and dig immediately into our discussion of black-hole astrophysics (\S\S{}3--8).
However, for the reader whose previous training has focussed primarily on
gravitation theory, the material of \S\ref{sec:2} is an important prerequisite for understanding
the rest of the notes.

The authors are deeply indebted to Mr.\ Alan Lightman for removing a large
number of errors from the manuscript.


% =============================================================================
% SECTION 2.1: Thermal Bremsstrahlung (pp.~347--348+)
% =============================================================================

\section{Thermal Bremsstrahlung (``Free-Free Radiation'' From a Plasma)}
\label{sec:2}

\subsection{Thermal Bremsstrahlung (``Free-Free Radiation'' From a Plasma)}
\label{sec:2.1}

\textit{Basic references}\quad chapter~6 of Shkarofsky, Johnston, and
Bachynski~(1966)\cite{Shkarofsky1966};
chapter~15 of Jackson~(1962)\cite{Jackson1962}; \S4 of
Ginzburg~(1967)\cite{Ginzburg1967}.

\bigskip

\textit{Notation for fundamental constants}

\medskip

\noindent
\begin{tabular}{ll}
$h$            & Planck's constant \\
$m_e$          & rest mass of electron \\
$m_p$          & rest mass of proton \\
$e$            & charge of electron \\
$c$            & speed of light \\
$\alpha$       & fine-structure constant \\
$r_0$          & classical electron radius \\
$R_y$          & Rydberg energy \\
$Z$            & charge of ions in plasma \\
$\mathcal{A}$  & atomic weight of ions \\
$k$            & Boltzmann's constant \\
$\xi$          & Euler constant ($= 1.78\ldots$)
\end{tabular}

\bigskip

\textit{Radiation from a single classical collision}

\medskip

Consider a single ion of charge~$Z$, at rest in the laboratory frame. Let a single
electron of speed $v \ll c$ (nonrelativistic) scatter off the nucleus (Coulomb
scattering), with impact parameter~$b$. Let $I(\omega)$ be the total energy per unit
circular frequency, $\omega$, radiated by the electron as it scatters. A simple classical
calculation [e.g.\ Jackson~(1962)\cite{Jackson1962}, p.~507] gives
%
\begin{equation}
I(\omega) = \frac{2}{3} \frac{e^2}{m_e c^3} \left|\Delta \mathbf{v}_\omega\right|^2,
\tag{2.1.1}
\end{equation}
%
where $\Delta \mathbf{v}_\omega$ is the change in electron velocity during a time $\tau = 1/\omega$ centered
about the electron's point of closest approach to the nucleus.

It is useful to examine two limiting cases: small-angle scattering ($\theta \ll 1$)
corresponding to large impact parameter ($b \gg b_{s\text{-}l}$); and large-angle scattering
($\theta \sim 2\pi$), corresponding to small impact parameter ($b \ll b_{s\text{-}l}$). The impact
parameter, $b_{s\text{-}l}$, which separates small-angle from large-angle scattering, is given
by
%
\begin{equation}
Ze^2 / b_{s\text{-}l} = \tfrac{1}{2}\, m_e v^2.
\tag{2.1.2}
\end{equation}

\textit{For small-angle scattering} the total scattering angle $\theta$ is
%
\begin{equation}
\theta = \frac{1}{v}\, |\Delta \mathbf{v}|
      = \frac{Ze^2/b}{\frac{1}{2}\, m_e v^2}
      = \frac{b_{s\text{-}l}}{b}\,,
\tag{2.1.3}
\end{equation}
%
and the time during which the scattering occurs is $\Delta t \approx b/v$.  Hence, for
$\tau = 1/\omega \gg \Delta t$, the change in velocity during time $\tau$, $|\Delta \mathbf{v}_\tau|$,
is the full change $|\Delta \mathbf{v}|$ of~(2.1.3); while for $\tau = 1/\omega \ll \Delta t$ the change is
essentially zero.  Inserting these changes into equation~(2.1.1), we obtain
%
\begin{align}
I(\omega) &= 0  &\text{for}&\quad \omega \gg v/b, \notag \\[4pt]
I(\omega) &= \frac{8}{3}\,\frac{Z^2 e^2}{m_e c}\left(\frac{r_0}{v\,b}\right)^{\!2}
           &\text{for}&\quad \omega \ll v/b.
\tag{2.1.4}
\end{align}

\textit{For large-angle scattering} the electron orbit is very nearly a parabola, with
``perihelion'' at radius
%
\begin{equation}
r_p = b^2 / b_{s\text{-}l}
\tag{2.1.5a}
\end{equation}
%
and with speed at perihelion
%
\begin{equation}
v_p = c\bigl(2 Z r_0 / r_p\bigr)^{1/2}.
\tag{2.1.5b}
\end{equation}
%
During time intervals $\tau = 1/\omega \ll r_p/v_p$, a negligible amount of velocity change
occurs; hence
%
\begin{equation}
I(\omega) = 0 \qquad \text{for}\quad \omega \gg v_p/r_p
           = c(2Zr_0\, b_{s\text{-}l}^{\,3})^{1/2}/b^3.
\tag{2.1.6a}
\end{equation}

During time intervals $\tau = 1/\omega \gg r_p/v_p$ but $\tau < b/v$, the electron is initially
($\tau = -\tau/2$) headed directly toward the nucleus with ``parabolic'' speed

% [Agent 01 translation ends here at bottom of book page 348.
%  Page 349 continues with eq.~(2.1.6b) and subsequent material.]
